\documentclass[11pt,a4paper]{article}
% Estilo comum dos PDFs individuais da Sprint 1.
\usepackage[utf8]{inputenc}
\usepackage[T1]{fontenc}
\usepackage[brazilian]{babel}
\usepackage[a4paper,margin=2.35cm]{geometry}
\usepackage{xcolor}
\usepackage{hyperref}
\usepackage{enumitem}
\usepackage{listings}
\usepackage{parskip}
\usepackage{needspace}
\definecolor{azul}{RGB}{20,63,112}
\hypersetup{colorlinks=true,linkcolor=azul,urlcolor=azul}
\setlist[itemize]{leftmargin=*,itemsep=3pt}
\lstset{basicstyle=\ttfamily\small,frame=single,breaklines=true,columns=fullflexible}
\sloppy
\newcommand{\codigo}[1]{\texttt{#1}}
\newcommand{\titulo}[2]{\begin{center}{\LARGE\bfseries #1\par}\vspace{4pt}{\large #2\par}\vspace{8pt}\rule{0.8\textwidth}{0.4pt}\par\vspace{5pt}{\small Matheus Moreli Durão Silva -- Matrícula 12621CBS217\par Cibersegurança - Bacharelado - Noturno\par}\end{center}}
\newcommand{\cardfonte}[1]{\par\smallskip\noindent\textit{Referência: }\href{#1}{Trello - enunciado e checklist do card}.}

\begin{document}\titulo{S1\_09 - Funções}{Reutilizar uma tarefa com parâmetros}
\section*{Entregável correto}
O arquivo deve conter uma função que recebe \codigo{nome} e \codigo{responsavel} e mostra o ativo formatado. O objetivo é perceber que uma função reúne uma tarefa com nome e pode receber dados sem depender de variáveis espalhadas pelo programa.
\section*{Leitura linha a linha}
\codigo{def} inicia a definição. Os nomes dentro dos parênteses são \textbf{parâmetros}: espaços reservados para os valores que virão na chamada. O bloco indentado é executado apenas quando a função é chamada. Neste exemplo, a função produz uma saída; quando há um valor para devolver, usa-se \codigo{return}, como em uma classificação de severidade.
\begin{lstlisting}
def mostrar_ativo(nome, responsavel):
    print(f"Ativo: {nome}")
    print(f"Responsavel: {responsavel}")

mostrar_ativo("Servidor de arquivos", "Equipe TI")
\end{lstlisting}
\section*{Como usar no projeto}
No \codigo{app.py}, uma função por responsabilidade torna o código menor de entender: \codigo{cadastrar\_ativo()}, \codigo{listar\_ativos()} e \codigo{salvar\_dados()} não misturam tarefas. Parâmetros permitem que uma função trabalhe com dados diferentes; \codigo{return} permite, por exemplo, devolver o ativo encontrado para outra parte do programa.
\section*{Checklist demonstrado}
Função, \codigo{def}, parâmetros, retorno e modularização foram estudados. Para conferir o entregável, trocar os argumentos da chamada e observar a mesma estrutura de saída com dados diferentes.
\cardfonte{https://trello.com/c/XWD40Pgo}
\end{document}
